% ============================================================================
% Section V: Software Architecture
% ============================================================================

\section{Software Architecture}
\label{sec:architecture}

\subsection{Design Principles}

The software architecture of \gtopt{} is organized around four core
principles:

\begin{enumerate}
  \item \textbf{Performance}: C++26 with cache-friendly sorted flat maps
    for sparse-matrix assembly, parallel assembly and solution of the
    (scene, phase) subproblems, and compressed in-memory LP snapshots.
  \item \textbf{Modularity}: Each power system component (generator,
    demand, line, battery, reservoir, etc.) has its own LP assembly class
    (\texttt{*\_lp.cpp}) that independently adds variables, constraints,
    and objective terms to the shared \texttt{LinearProblem}.
  \item \textbf{Extensibility}: New components and solver backends can be
    added without modifying the core LP engine.  Solver backends are
    loaded as dynamic plugins (\texttt{libgtopt\_solver\_*.so}).
  \item \textbf{Interoperability}: JSON configuration with native Apache
    Arrow/Parquet I/O enables seamless integration with Python, R, and
    other data science ecosystems.
\end{enumerate}

\subsection{System Architecture}

Fig.~\ref{fig:arch_layers} illustrates the layered architecture.

\begin{figure}[t]
  \centering
  \begin{tikzpicture}[
    >=Stealth,
    layer/.style={draw, thick, rounded corners=3pt, minimum width=7cm,
      minimum height=1.1cm, font=\small, align=center},
    sublbl/.style={font=\tiny, text=black!60},
    arr/.style={->, thick, gray!70},
  ]
    % Layer 4 (bottom): Data Layer
    \node[layer, fill=green!12] (data) at (0,0)
      {\textbf{Data Layer}};
    \node[sublbl, below=-0.15cm of data]
      {JSON \quad Parquet/CSV \quad Schedule Resolution};

    % Layer 3: Model Layer
    \node[layer, fill=blue!10, above=0.5cm of data] (model)
      {\textbf{Model Layer}};
    \node[sublbl, below=-0.15cm of model]
      {Planning $\to$ System $\to$ Components};

    % Layer 2: LP Assembly Layer
    \node[layer, fill=orange!12, above=0.5cm of model] (lp)
      {\textbf{LP Assembly Layer}};
    \node[sublbl, below=-0.15cm of lp]
      {PlanningLP $\to$ SimulationLP $\to$ Component LPs};

    % Layer 1 (top): Solver Layer
    \node[layer, fill=red!10, above=0.5cm of lp] (solver)
      {\textbf{Solver Layer}};
    \node[sublbl, below=-0.15cm of solver]
      {HiGHS \quad CBC/CLP \quad SCIP \quad cuOpt \quad CPLEX \quad Gurobi \quad MindOpt};

    % Arrows between layers
    \draw[arr] (data.north) -- (model.south);
    \draw[arr] (model.north) -- (lp.south);
    \draw[arr] (lp.north) -- (solver.south);

    % Side labels
    \node[font=\tiny, rotate=90, anchor=south, gray]
      at (-3.9,1.2) {Data flow};
    \draw[->, thin, gray] (-3.9,0) -- (-3.9,3.6);
  \end{tikzpicture}
  \caption{\gtopt{} layered architecture: data ingestion, model
    construction, LP assembly, and pluggable solver backends.}
  \label{fig:arch_layers}
\end{figure}

The software is organized into four layers:

\begin{enumerate}
  \item \textbf{Data Layer}: JSON parsing (via DAW JSON Link, a
    zero-allocation header-only library), Parquet/CSV I/O (via Apache
    Arrow C++ SDK), and schedule resolution (\texttt{FieldSched<T>}
    variant type supporting scalar, vector, or filename references).
  \item \textbf{Model Layer}: The \texttt{Planning} $\rightarrow$
    \texttt{System} $\rightarrow$ component hierarchy that maps JSON
    input to C++ data structures.
  \item \textbf{LP Assembly Layer}: \texttt{PlanningLP} orchestrates
    \texttt{SimulationLP} instances (one per scenario/stage), which in
    turn assemble component LP objects (\texttt{GeneratorLP},
    \texttt{DemandLP}, \texttt{LineLP}, \texttt{BatteryLP},
    \texttt{ReservoirLP}, etc.) into a shared
    \texttt{FlatLinearProblem}.
  \item \textbf{Solver Layer}: The \texttt{SolverBackend} interface
    abstracts LP/MIP solution.  Implementations are loaded as dynamic
    plugins from configurable search paths.
\end{enumerate}

\subsection{Data Model and JSON Interface}

The input is specified as one or more JSON files containing three sections:

\begin{lstlisting}[language={},caption={Top-level JSON structure}]
{
  "options": { ... },
  "simulation": {
    "block_array": [...],
    "stage_array": [...],
    "scenario_array": [...]
  },
  "system": {
    "bus_array": [...],
    "generator_array": [...],
    "demand_array": [...],
    "line_array": [...],
    "battery_array": [...],
    ...
  }
}
\end{lstlisting}

Time-series data (demand profiles, generator capacity factors, reservoir
inflows) can be specified inline as JSON arrays, or as references to
external Parquet or CSV files in the \texttt{input\_directory}.  This
hybrid approach keeps the JSON configuration compact while supporting
arbitrarily large time-series via columnar Parquet storage.  Elements
can be referenced by name or by numeric identifier.  The command
\texttt{gtopt --json-schema} prints the JSON Schema of the input,
generated from the C++ types, for editor completion and validation.
Results are written per (scene, phase) as Parquet files (snappy or zstd)
or CSV, readable by pandas, DuckDB, Spark or Power BI.

\subsection{User Constraints: PAMPL}

Operating rules that are not part of the element library are written in
PAMPL, an AMPL-like language read from a text file.  PAMPL constraints
reference element variables by name, use parameters, sums over element
sets and scenario/stage/block domains, and can include state terms that
cross phases in SDDP.  The Chilean irrigation agreements of the Maule
and Laja basins, for example, are built from water-right elements
(Section~\ref{sec:formulation}) plus PAMPL constraints, without changes
to the C++ code.

\subsection{Auditability}

Every LP/MIP can be written as a standard LP file with readable row and
column names, and solved or checked with any external solver.  When a
run is infeasible, the LP is dumped and \texttt{gtopt\_check\_lp}
computes an irreducible infeasible subsystem.  An LP
\emph{fingerprint} records the structure of the model actually solved
(element classes, constraint families and their counts), so that
auditors and continuous integration can verify that a run uses the
intended formulation.  The complete formulation, including the SDDP
cut-validity conditions, is published with the source code.

\subsection{Python Toolchain}

A Python toolchain is provided for data preparation, conversion and
validation:

\begin{itemize}
  \item \textbf{igtopt}: Excel workbook ($\rightarrow$ JSON + Parquet)
    converter for spreadsheet-based case preparation.
  \item \textbf{Converters} from the models used in the region and in
    the literature: \texttt{plp2gtopt} (PLP, including its cuts),
    \texttt{plexos2gtopt} (PLEXOS inputs), \texttt{sddp2gtopt}
    (PSR SDDP), PSS/E, NREL Sienna, RTS-GMLC and NREL-118.
  \item \textbf{pp2gtopt / gtopt2pp}: pandapower
    ($\leftrightarrow$ gtopt) bidirectional converters for power flow
    model exchange.
  \item \textbf{gtopt\_compare}: Automated cross-validation against
    pandapower DC OPF on IEEE test networks.
  \item \textbf{gtopt\_diagram}: network and hydro-topology diagrams.
  \item \textbf{gtopt-net}: network reduction for long-term studies,
    with decisions mapped back onto the full network.
  \item \textbf{run\_gtopt / gtopt\_monitor}: execution management with
    a terminal dashboard and live SDDP convergence monitoring.
  \item \textbf{Validation tools}: \texttt{gtopt\_check\_json},
    \texttt{gtopt\_check\_lp}, \texttt{gtopt\_check\_output},
    \texttt{gtopt\_check\_solvers} for input validation and diagnostics.
\end{itemize}

\subsection{Web and GUI Services}

Two companion services provide graphical interfaces:

\begin{itemize}
  \item \textbf{webservice}: a Next.js REST API and browser UI for
    submitting optimization jobs, monitoring progress and browsing
    results.
  \item \textbf{guiservice}: a Python/Flask application for creating
    and editing cases, viewing the network topology and plotting
    results in the browser.
\end{itemize}

\subsection{Engineering Practice}

The C++ library, the Python tools and the web services are covered by
unit and integration test suites (more than 5\,000 C++ test cases),
static analysis (clang-tidy and CodeQL) and coverage measurement, all
run in continuous integration on every change.  The test suite includes
the benchmark cases of Section~\ref{sec:validation}, reference results
for each case, and solver-specific tolerances where backends may return
different optimal dual vertices.
